\begin{lemma}
\label{lemma-alternating-cech-complex-complex}
In Situation \ref{situation-complex}. Let $M^\bullet$ be a
complex of $A$-modules and
denote $\mathcal{F}^\bullet$ the associated complex of
$\mathcal{O}_X$-modules. Then
there is a canonical isomorphism of complexes
$$
\colim_e \Hom^\bullet(I^\bullet(f_1^e, \ldots, f_r^e), M^\bullet)
\longrightarrow
\text{Tot}(\check{\mathcal{C}}_{alt}^\bullet(\mathcal{U}, \mathcal{F}^\bullet))
$$
functorial in $M^\bullet$.
\end{lemma}

\begin{proof}
Consider the double complex $F^{\bullet, \bullet}$ with terms
$F^{p, q} = \mathcal{C}_{alt}^p(\mathcal{U}, \mathcal{F}^q)$
discussed in
Cohomology, Section \ref{cohomology-section-cech-cohomology-of-complexes}.
Consider the double complex $G^{\bullet, \bullet}$ with terms
$G^{p, q} = \colim_e \Hom_A(I^{-p}(f_1^e, \ldots, f_r^e), M^q)$ and
differentials given by functoriality (without the intervention
of signs). The maps $\psi^{p, q} : G^{p, q} \to F^{p, q}$
constructed in the proof of Lemma \ref{lemma-alternating-cech-complex}
are isomorphisms and compatible with the differentials $d_1$ (by the lemma)
and $d_2$ (this is clear). However, the differentials $d$ on the
complexes on the left and right hand side of the arrow in the lemma
have different signs. Namely, for $g \in G^{p, q}$ is given by
$$
d(g) =  d_2(g) - (-1)^{p + q} d_1(g) 
$$
(see More on Algebra, Section \ref{more-algebra-section-hom-complexes})
and the differential for $f \in F^{p, q}$ is given by
$$
d(f) = d_1(f) + (-1)^p d_2(f)
$$
Thus we can fix the signs by multiplying $\psi^{p, q}$ by
$(-1)^{pq + p(p - 1)/2}$.
\end{proof}
